% id: 71-07
% book: 베이직쎈 미적분1 · 04 도함수의 활용 (1)
% section: 유형 045 평균값 정리
% figure: none
% answer: 
% answer_source: 
% variant_level: 0 (원본 전사)
% 이 파일은 build.mjs 가 items.json 에서 만든다. 고칠 때는 items.json 을 고친다.
\smash{\raisebox{1.5mm}{\dmkichul{베이직쎈 미적분1 71쪽 07번}}}
다음은 함수 $f(x)$가 닫힌구간 $[a{,}\mkern3mu\mathopen{}b]$에서 연속이고 열린구간 $(a{,}\mkern3mu\mathopen{}b)$에서 미분가능할 때, 열린구간 $(a{,}\mkern3mu\mathopen{}b)$에 속하는 모든 $x$에 대하여 $f'(x)=0$이면 함수 $f(x)$는 닫힌구간 $[a{,}\mkern3mu\mathopen{}b]$에서 상수함수임을 증명하는 과정이다. ㈎, ㈏에 알맞은 것을 구하시오. \exprbox{\begin{minipage}{0.96\linewidth}\raggedright\lineskip=4pt {\small\bfseries 증명}\par\smallskip $a<x\le b$인 임의의 실수 $x$에 대하여 함수 $f(x)$가 닫힌구간 $[a{,}\mkern3mu\mathopen{}x]$에서 연속이고 열린구간 $(a{,}\mkern3mu\mathopen{}x)$에서 미분가능하므로 평균값 정리에 의하여\par\smallskip\dispeq{$\dfrac{f(x)-f(a)}{x-a}=f'(c)$}\par\smallskip\noindent 인 $c$가 열린구간 \framebox[2.2em]{\scriptsize ㈎}에 적어도 하나 존재한다.\\ 그런데 $f'(c)=0$이므로 \quad $f(x)-\framebox[2.2em]{\scriptsize ㈏}=0$\\ \hspace*{1.5em}$\therefore\ f(x)=\framebox[2.2em]{\scriptsize ㈏}$\\ 따라서 함수 $f(x)$는 닫힌구간 $[a{,}\mkern3mu\mathopen{}b]$에서 상수함수이다.\end{minipage}}
